\section{Iterative estimates for eigenvalues}
Follow the book's indexing carefully. Starting with $x_0$, compute
$y_k=Ax_k$, choose an entry $\mu_k$ of $y_k$ with largest absolute
value (retain its sign), and set $x_{k+1}=y_k/\mu_k$.
Thus $\mu_1$ comes from $Ax_1$, not $Ax_0$. A negative dominant
eigenvalue need not cause the normalized vectors to alternate, because
the signed normalization removes the reversal.

For inverse iteration with a shift $\alpha$, solve
$(A-\alpha I)y_k=x_k$ instead. Normalize by its largest-magnitude
entry $\mu_k$, and estimate the original eigenvalue by
$\theta_k=\alpha+1/\mu_k$. Reuse a factorization of $A-\alpha I$;
do not explicitly form its inverse at every step.
A small residual $\|Av-\theta v\|$ checks a proposed eigenpair.
Agreement of two consecutive estimates alone can be misleading.

\practice{5.8.P1}{Test proportionality and inspect the residual}
The given vector $x=(1,-4.3,8.1)^T$ gives $Ax=(3,-13,24.5)^T$.
The entrywise ratios are $3$, $130/43\approx3.023256$, and
$245/81\approx3.024691$. Their close agreement suggests an eigenvalue
near $3.024$. The Rayleigh estimate is
\[R(x)=\frac{x^TAx}{x^Tx}=\frac{25735}{8510}
\approx3.024089.\]
For $\theta=R(x)$, inspect $Ax-\theta x$ relative to the size of
$Ax$ or $x$. A residual near zero, rather than the raw size of $x$,
is the meaningful approximation test. Exact proportionality would give
zero residual and an exact eigenvector.

\exercise{5.8.1}{Read the limiting direction and retain the scaling sign}
The final supplied normalized vector is approximately $\left[\begin{matrix}1\\\frac{1663}{5000}\end{matrix}\right]$.
Multiplication by the stated matrix gives
\[Ax\approx\left[\begin{matrix}\frac{24989}{5000}\\\frac{4163}{2500}\end{matrix}\right].\]
The largest-magnitude component is $\mu\approx4.997800$.
The trend in the vectors suggests the sharper limiting direction
$v=\left[\begin{matrix}1\\\frac{1}{3}\end{matrix}\right]$, and direct multiplication checks
$Av=5v$ exactly.
\answer{The dominant eigenvalue is $5$, with eigenvector
$\left[\begin{matrix}1\\\frac{1}{3}\end{matrix}\right]$.}

\exercise{5.8.2}{Read the limiting direction and retain the scaling sign}
The final supplied normalized vector is approximately $\left[\begin{matrix}- \frac{63}{250}\\1\end{matrix}\right]$.
Multiplication by the stated matrix gives
\[Ax\approx\left[\begin{matrix}- \frac{1567}{1250}\\\frac{3129}{625}\end{matrix}\right].\]
The largest-magnitude component is $\mu\approx5.006400$.
The trend in the vectors suggests the sharper limiting direction
$v=\left[\begin{matrix}- \frac{1}{4}\\1\end{matrix}\right]$, and direct multiplication checks
$Av=5v$ exactly.
\answer{The dominant eigenvalue is $5$, with eigenvector
$\left[\begin{matrix}- \frac{1}{4}\\1\end{matrix}\right]$.}

\exercise{5.8.3}{Read the limiting direction and retain the scaling sign}
The final supplied normalized vector is approximately $\left[\begin{matrix}\frac{1297}{2500}\\1\end{matrix}\right]$.
Multiplication by the stated matrix gives
\[Ax\approx\left[\begin{matrix}\frac{2297}{5000}\\\frac{2836}{3125}\end{matrix}\right].\]
The largest-magnitude component is $\mu\approx0.907520$.
The trend in the vectors suggests the sharper limiting direction
$v=\left[\begin{matrix}\frac{1}{2}\\1\end{matrix}\right]$, and direct multiplication checks
$Av=\frac{9}{10}v$ exactly.
\answer{The dominant eigenvalue is $\frac{9}{10}$, with eigenvector
$\left[\begin{matrix}\frac{1}{2}\\1\end{matrix}\right]$.}

\exercise{5.8.4}{Read the limiting direction and retain the scaling sign}
The final supplied normalized vector is approximately $\left[\begin{matrix}1\\\frac{3751}{5000}\end{matrix}\right]$.
Multiplication by the stated matrix gives
\[Ax\approx\left[\begin{matrix}- \frac{1003}{2500}\\- \frac{3761}{12500}\end{matrix}\right].\]
The largest-magnitude component is $\mu\approx-0.401200$.
The trend in the vectors suggests the sharper limiting direction
$v=\left[\begin{matrix}1\\\frac{3}{4}\end{matrix}\right]$, and direct multiplication checks
$Av=- \frac{2}{5}v$ exactly.
\answer{The dominant eigenvalue is $- \frac{2}{5}$, with eigenvector
$\left[\begin{matrix}1\\\frac{3}{4}\end{matrix}\right]$.} Here ``largest'' in the exercise means
largest in \emph{absolute value}: the other eigenvalue is $.1$, which
is algebraically larger but not dominant.

\exercise{5.8.5}{Normalize by the second component of the supplied iterate}
Dividing the last vector by its second entry gives
\[x=\left[\begin{matrix}- \frac{24991}{31241}\\1\end{matrix}\right]\approx\left[\begin{matrix}- \frac{7999}{10000}\\1\end{matrix}\right].\]
Use the requested four-decimal vector for a check:
\[A\left[\begin{matrix}- \frac{7999}{10000}\\1\end{matrix}\right]=\left[\begin{matrix}\frac{8003}{2000}\\- \frac{2501}{500}\end{matrix}\right].\]
The second entry estimates the eigenvalue by $- \frac{2501}{500}$,
since the corresponding input entry is $1$.
The nearby exact eigenpair is $\lambda=-5$,
$v=\left[\begin{matrix}- \frac{4}{5}\\1\end{matrix}\right]$, verified by direct multiplication.
The rounded estimate is approximate, so exact equality $Ax=\lambda x$
is not expected until using the exact vector.

\exercise{5.8.6}{Normalize by the second component of the supplied iterate}
Dividing the last vector by its second entry gives
\[x=\left[\begin{matrix}- \frac{2045}{4093}\\1\end{matrix}\right]\approx\left[\begin{matrix}- \frac{1249}{2500}\\1\end{matrix}\right].\]
Use the requested four-decimal vector for a check:
\[A\left[\begin{matrix}- \frac{1249}{2500}\\1\end{matrix}\right]=\left[\begin{matrix}- \frac{2501}{1250}\\\frac{5003}{1250}\end{matrix}\right].\]
The second entry estimates the eigenvalue by $\frac{5003}{1250}$,
since the corresponding input entry is $1$.
The nearby exact eigenpair is $\lambda=4$,
$v=\left[\begin{matrix}- \frac{1}{2}\\1\end{matrix}\right]$, verified by direct multiplication.
The rounded estimate is approximate, so exact equality $Ax=\lambda x$
is not expected until using the exact vector.

\exercise{5.8.7}{Compute the requested five indexed power-method rows}
Use the signed largest-magnitude entry for normalization and do not
round intermediate vectors. For $A=\left[\begin{matrix}6 & 7\\8 & 5\end{matrix}\right]$ and $x_0=(1,0)^T$,
the results are
\[\begin{array}{c|rr|r}k&(x_k)_1&(x_k)_2&\mu_k\\\hline
1&0.750000000&1.000000000&11.500000000\\
2&1.000000000&0.956521739&12.782608696\\
3&0.993197279&1.000000000&12.959183673\\
4&1.000000000&0.998950131&12.994750656\\
5&0.999838416&1.000000000&12.999030499\end{array}.\]
For an indexing check, the first multiplication gives
$Ax_0=\left[\begin{matrix}6\\8\end{matrix}\right]$ and
$\mu_0=8$. Dividing by this scalar produces
$x_1$, after which $\mu_1$ is computed from $Ax_1$.
The values converge toward the dominant eigenvalue $13$,
with eigenvector direction $(1,1)^T$.

\exercise{5.8.8}{Compute the requested five indexed power-method rows}
Use the signed largest-magnitude entry for normalization and do not
round intermediate vectors. For $A=\left[\begin{matrix}2 & 1\\4 & 5\end{matrix}\right]$ and $x_0=(1,0)^T$,
the results are
\[\begin{array}{c|rr|r}k&(x_k)_1&(x_k)_2&\mu_k\\\hline
1&0.500000000&1.000000000&7.000000000\\
2&0.285714286&1.000000000&6.142857143\\
3&0.255813953&1.000000000&6.023255814\\
4&0.250965251&1.000000000&6.003861004\\
5&0.250160772&1.000000000&6.000643087\end{array}.\]
For an indexing check, the first multiplication gives
$Ax_0=\left[\begin{matrix}2\\4\end{matrix}\right]$ and
$\mu_0=4$. Dividing by this scalar produces
$x_1$, after which $\mu_1$ is computed from $Ax_1$.
The values converge toward the dominant eigenvalue $6$,
with eigenvector direction $(1/4,1)^T$.

\exercise{5.8.9}{Continue the same indexed algorithm to the sixth estimate}
For the given matrix
\[A=\left[\begin{matrix}8 & 0 & 12\\1 & -2 & 1\\0 & 3 & 0\end{matrix}\right],\qquad x_0=(1,0,0)^T,\]
repeatedly compute $Ax_k$, retain its signed largest-magnitude entry,
and divide to form $x_{k+1}$. This gives
\[\mu_5\approx8.4233109850,\qquad
\mu_6\approx8.4246301920.\]
The accompanying current vectors, which make these estimates reproducible,
are
\[x_5\approx\left[\begin{matrix}1.0\\0.0993553378\\0.0352759154\end{matrix}\right],\qquad
x_6\approx\left[\begin{matrix}1.0\\0.0993154879\\0.0353858493\end{matrix}\right].\]
Multiplying each vector by $A$ and selecting its largest-magnitude
entry checks the listed scalars. The vectors have been rounded only
for display, so the check agrees to the displayed precision.

\exercise{5.8.10}{Continue the same indexed algorithm to the sixth estimate}
For the given matrix
\[A=\left[\begin{matrix}1 & 2 & -2\\1 & 1 & 9\\0 & 1 & 9\end{matrix}\right],\qquad x_0=(1,0,0)^T,\]
repeatedly compute $Ax_k$, retain its signed largest-magnitude entry,
and divide to form $x_{k+1}$. This gives
\[\mu_5\approx9.9319371728,\qquad
\mu_6\approx9.9871727289.\]
The accompanying current vectors, which make these estimates reproducible,
are
\[x_5\approx\left[\begin{matrix}0.0183246073\\1.0\\0.9904013962\end{matrix}\right],\qquad
x_6\approx\left[\begin{matrix}0.0037778949\\1.0\\0.9981549815\end{matrix}\right].\]
Multiplying each vector by $A$ and selecting its largest-magnitude
entry checks the listed scalars. The vectors have been rounded only
for display, so the check agrees to the displayed precision.

\exercise{5.8.11}{Compare the scaling estimate with the Rayleigh quotient}
The exact dominant eigenvalue, checked from the characteristic
polynomial, is $6$. Compute each Rayleigh quotient
from the current $x_k$ using $R(x_k)=x_k^TAx_k/(x_k^Tx_k)$:
\[\begin{array}{c|rr|rr}
k&\mu_k&R(x_k)&|\mu_k-\lambda|&|R(x_k)-\lambda|\\\hline
1&5.800000000000&5.965517241379&2.000e-01&3.448e-02\\
2&5.965517241379&5.999035679846&3.448e-02&9.643e-04\\
3&5.994219653179&5.999973208305&5.780e-03&2.679e-05\\
4&5.999035679846&5.999999255782&9.643e-04&7.442e-07\end{array}.\]
The quotient's error is substantially smaller at the same iterate.
Both matrices are symmetric, so this demonstrates the improved
accuracy claimed in the exercise. Signed normalization is especially
important for the negative dominant eigenvalue in Exercise 12.

\exercise{5.8.12}{Compare the scaling estimate with the Rayleigh quotient}
The exact dominant eigenvalue, checked from the characteristic
polynomial, is $-4$. Compute each Rayleigh quotient
from the current $x_k$ using $R(x_k)=x_k^TAx_k/(x_k^Tx_k)$:
\[\begin{array}{c|rr|rr}
k&\mu_k&R(x_k)&|\mu_k-\lambda|&|R(x_k)-\lambda|\\\hline
1&-4.333333333333&-3.923076923077&3.333e-01&7.692e-02\\
2&-3.923076923077&-3.995121951220&7.692e-02&4.878e-03\\
3&-4.019607843137&-3.999694842844&1.961e-02&3.052e-04\\
4&-3.995121951220&-3.999980926586&4.878e-03&1.907e-05\end{array}.\]
The quotient's error is substantially smaller at the same iterate.
Both matrices are symmetric, so this demonstrates the improved
accuracy claimed in the exercise. Signed normalization is especially
important for the negative dominant eigenvalue in Exercise 12.

\exercise{5.8.13}{A theoretically convergent method can be impractically slow}
If the two leading eigenvalues have different moduli, there is a
strictly dominant one, so the power method converges for an initial
vector with a nonzero component in that eigendirection (under the
section's eigenbasis hypotheses). But the ratio of the other large
modulus to the dominant modulus is near $4/4=1$. The unwanted
mode therefore decays slowly relative to the desired one.
\answer{It can work, but may require many iterations and be less
useful than shifted inverse iteration. An initial vector with no
dominant component is an exception to the usual convergence claim.}

\exercise{5.8.14}{Break the tie by shifting toward the positive eigenvalue}
If the leading eigenvalues are exactly $\lambda$ and $-\lambda$,
their powers have equal moduli; ordinary power iteration need not
select a single direction. Choose a shift $\alpha$ near the positive
eigenvalue (for example $4.1$, if suitably close) but not exactly an
eigenvalue. Inverse iteration applies to $(A-\alpha I)^{-1}$.
Its eigenvalues are $1/(\lambda_j-\alpha)$, so the desired positive
eigenvalue is dominant when its distance from the shift is uniquely
smallest. Solving shifted systems and normalizing then gives a
sequence for that positive eigenpair.
Alternatively, adding a suitable positive multiple of $I$ can create
a unique positive dominant modulus; subtract that shift from the
resulting eigenvalue estimate. The inverse method is usually faster
when an accurate initial estimate is available.

\exercise{5.8.15}{Transform an eigenvalue under a shifted inverse}
From $Ax=\lambda x$ subtract $\alpha x$ to get
$(A-\alpha I)x=(\lambda-\alpha)x$.
Since $\alpha$ is not an eigenvalue, the matrix is invertible and
$\lambda-\alpha\ne0$. Multiply by $(A-\alpha I)^{-1}$:
\[x=(\lambda-\alpha)(A-\alpha I)^{-1}x,
\qquad (A-\alpha I)^{-1}x=\frac1{\lambda-\alpha}x.\]
Because $x\ne0$, this is exactly the claimed eigenpair.

\exercise{5.8.16}{Recover the original eigenvalue from the inverse one}
Multiply $(A-\alpha I)^{-1}x=\mu x$ by $A-\alpha I$:
$x=\mu(A-\alpha I)x$.
The inverse is invertible, so its eigenvalue $\mu$ is nonzero.
Dividing by $\mu$ and rearranging gives
\[Ax=\left(\alpha+\frac1\mu\right)x.\]
Thus $\alpha+1/\mu$ is an eigenvalue of $A$, with the same
nonzero eigenvector. This proves the scalar update in inverse iteration.

\exercise{5.8.17}{Run shifted inverse iteration and verify its residual}
Use
\[A=\left[\begin{matrix}10 & -8 & -4\\-8 & 13 & 4\\-4 & 5 & 4\end{matrix}\right],\quad \alpha=3.3,\quad x_0=(1,0,0)^T.\]
At each step solve $(A-\alpha I)y_k=x_k$, choose a signed
largest-magnitude component $\mu_k$, and set
$x_{k+1}=y_k/\mu_k$, $\theta_k=\alpha+1/\mu_k$.
Using high-precision arithmetic, iteration index $k=7$ gives
\[\theta_k\approx3.3212201246571,\qquad
x_{k+1}\approx\left[\begin{matrix}1.0\\0.787012129501\\0.0956707098337\end{matrix}\right].\]
Before rounding, $\|Ax_{k+1}-\theta_kx_{k+1}\|_\infty<10^{-10}$
and the last two scalar estimates differ by less than $10^{-12}$.
The eigenvalue to four decimal places is $\boxed{3.3212}$.
Independent roots of the exact characteristic polynomial confirm this
choice. The shift makes this eigenvalue uniquely closest, so inverse
iteration selects it rather than the dominant eigenvalue of $A$.

\exercise{5.8.18}{Run shifted inverse iteration and verify its residual}
Use
\[A=\left[\begin{matrix}8 & 0 & 12\\1 & -2 & 1\\0 & 3 & 0\end{matrix}\right],\quad \alpha=-1.4,\quad x_0=(1,0,0)^T.\]
At each step solve $(A-\alpha I)y_k=x_k$, choose a signed
largest-magnitude component $\mu_k$, and set
$x_{k+1}=y_k/\mu_k$, $\theta_k=\alpha+1/\mu_k$.
Using high-precision arithmetic, iteration index $k=10$ gives
\[\theta_k\approx-1.424428900898,\qquad
x_{k+1}\approx\left[\begin{matrix}1.0\\0.372900802803\\-0.785369075075\end{matrix}\right].\]
Before rounding, $\|Ax_{k+1}-\theta_kx_{k+1}\|_\infty<10^{-10}$
and the last two scalar estimates differ by less than $10^{-12}$.
The eigenvalue to four decimal places is $\boxed{-1.4244}$.
Independent roots of the exact characteristic polynomial confirm this
choice. The shift makes this eigenvalue uniquely closest, so inverse
iteration selects it rather than the dominant eigenvalue of $A$.

\exercise{5.8.19}{Compute both extreme-modulus eigenpairs from the required start}
Use $x_0=(1,0,0,0)^T$ for both computations.
\step{(a) Ordinary power iteration.}
With signed largest-component normalization, the scalar estimate
converges to
\[\lambda_{\max}\approx30.288685345802,\quad
v_{\max}\approx\left[\begin{matrix}0.95762878047\\0.68893692056\\1.0\\0.94378150874\end{matrix}\right].\]
The four-decimal eigenvalue is $\boxed{30.2887}$.
\step{(b) Inverse iteration with shift zero.}
Repeatedly solve $Ay_k=x_k$ instead of explicitly multiplying by
$A^{-1}$. The smallest-modulus eigenpair is
\[\lambda_{\mathrm{near}\,0}\approx0.010150048397894,\quad
v_{\mathrm{near}\,0}\approx\left[\begin{matrix}-0.60397234234\\1.0\\-0.25113513051\\0.14895344556\end{matrix}\right].\]
The four-decimal eigenvalue is $\boxed{0.0102}$.
The stopping tests were both a scalar-estimate change below $10^{-12}$
and an unrounded residual below $10^{-10}$. They hold at indices
$k=16$ and $k=6$, respectively, with $k=0$
representing the first multiplication or solve.
Independent characteristic-polynomial roots verify which eigenvalues
are largest and nearest zero. The displayed eigenvectors are scaled
to make their largest-magnitude component equal to $1$.

\exercise{5.8.20}{Compute both extreme-modulus eigenpairs from the required start}
Use $x_0=(1,0,0,0)^T$ for both computations.
\step{(a) Ordinary power iteration.}
With signed largest-component normalization, the scalar estimate
converges to
\[\lambda_{\max}\approx19.182036763332,\quad
v_{\max}\approx\left[\begin{matrix}0.1844163575\\1.0\\0.17961525893\\0.40710960748\end{matrix}\right].\]
The four-decimal eigenvalue is $\boxed{19.1820}$.
\step{(b) Inverse iteration with shift zero.}
Repeatedly solve $Ay_k=x_k$ instead of explicitly multiplying by
$A^{-1}$. The smallest-modulus eigenpair is
\[\lambda_{\mathrm{near}\,0}\approx0.012205562828844,\quad
v_{\mathrm{near}\,0}\approx\left[\begin{matrix}1.0\\0.22260978915\\-0.91799334087\\0.66048300357\end{matrix}\right].\]
The four-decimal eigenvalue is $\boxed{0.0122}$.
The stopping tests were both a scalar-estimate change below $10^{-12}$
and an unrounded residual below $10^{-10}$. They hold at indices
$k=18$ and $k=6$, respectively, with $k=0$
representing the first multiplication or solve.
Independent characteristic-polynomial roots verify which eigenvalues
are largest and nearest zero. The displayed eigenvectors are scaled
to make their largest-magnitude component equal to $1$.

\exercise{5.8.21}{Separate vector size from convergence of direction}
For $x=(.5,.5)^T$, diagonal multiplication is explicit:
\begin{enumerate}[label=(\alph*)]
\item $A^kx=(.5(.8)^k,.5(.2)^k)^T\to0$.
The zero vector is not an eigenvector. Yet the second coordinate
relative to the first is $(1/4)^k\to0$, so the direction tends to
the first coordinate axis.
\item $A^kx=(.5,.5(.8)^k)^T\to(.5,0)^T$, a nonzero eigenvector.
Here the dominant eigenvalue is exactly $1$, so its scale is preserved.
\item $A^kx=(.5(8)^k,.5(2)^k)^T$ grows without bound.
Its direction again tends to the first coordinate axis, but it does
not converge to a finite vector.
\end{enumerate}
The common conclusion is about normalized direction, not convergence
of $A^kx$ itself. Moreover, a vector exactly missing the dominant mode
cannot reveal that eigendirection. Appropriate normalization is therefore
an essential part of the power method, not an optional cosmetic step.

\textbf{Reproducing the iterative computations.}
The following complete functions implement the signed indexing used
in the tables. Enter the matrix printed in the chosen exercise and its
stated initial vector as NumPy arrays. For 19--20 and inverse iteration,
continue until both the estimate change and eigenpair residual are small.
\begin{verbatim}
import numpy as np
def power_rows(A, x, count):
    for k in range(count+1):
        y = A @ x
        mu = y[np.argmax(np.abs(y))]
        rayleigh = (x @ A @ x)/(x @ x)
        yield k, x.copy(), mu, rayleigh
        x = y/mu

def inverse_rows(A, alpha, x, count):
    H = A-alpha*np.eye(len(x))
    for k in range(count+1):
        y = np.linalg.solve(H, x)
        mu = y[np.argmax(np.abs(y))]
        xnew = y/mu
        theta = alpha+1/mu
        residual = np.max(np.abs(A@xnew-theta*xnew))
        yield k, xnew.copy(), theta, residual
        x = xnew
\end{verbatim}
All numerical results above were actually computed from the supplied
matrices; the code is provided to reproduce them. Extra working precision
was used for the final residual checks in Exercises 17--20.
